\begin{textsample}{POS Dim 6 – vem\_ed\_06 – Score 2.00 – vem\_ed\_06/t021.txt}  \label{ex:f6_pos_076}
VEm Virtual Exchange Medium QUEBRA-GELO Conforme pesquisa realizada pela equipe dos PCIs no segundo semestre de 2020, 75 \% dos Projetos Colaborativos Internacionais das Fatecs são desenvolvidos em língua inglesa e 25 \% em espanhol.
Uma nova frente de trabalho se abre, para intercâmbios virtuais realizados em língua \textbf{portuguesa}, a partir das conversas realizadas com o Instituto Politécnico de Viseu, em Portugal.
Assim, surgem oportunidades para \textbf{aumentar} a proficiência dos estudantes no próprio idioma nativo, além de desenvolver competências interculturais, digitais e de trabalho em equipe internacional.
Na página 3, você fica conhecendo os primeiros passos de uma colaboração que promete bons frutos.
Eu, que sempre realizo reuniões em inglês ou espanhol, fico emocionado em falar \textbf{português} com os colegas de Viseu!
É um aprendizado constante.
Na seção “Quem é quem” (página 2), uma breve entrevista com Lavern Samuels, diretor do escritório internacional da Durban University of Technology (DUT, África do Sul).
Ele comenta sobre a importância das colaborações Sul-Sul.
Em “Boas Práticas”, na página 4, Neusa Haruka Sezaki Gritti, da nossa equipe dos PCIs, compartilha a \textbf{evolução} do projeto desenvolvido entre Tianjin Normal University, na China, e as Fatecs de Franca, Campinas, Catanduva, São José do Rio Preto, Itaquaquecetuba, Piracicaba, Sebrae (na capital) e Mogi das Cruzes.
Ao longo de cinco edições, a colaboração amadureceu, focada na aprendizagem de inglês por falantes não nativos, com trocas culturais riquíssimas sobre jeitos de ser nos dois países.
As dicas de Neusa para o sucesso de um PCI são preciosas para professores interessados em desenvolver intercâmbios virtuais, independentemente do tema a ser tratado no projeto.
Boa leitura!

% matched lemmas: aumentar, evolução, português
\end{textsample}
